


As the complexity of datasets has grown in recent years, topological methods have emerged as powerful complements to state-of-the-art ML methods. The advent of ML has revolutionized the way we analyze and interpret complex data, yet there remain challenges in capturing the intrinsic topological structures inherent in such data. Traditional ML techniques, while powerful, often fall short in identifying and leveraging these structures, leading to the potential loss of valuable insights. \textit{Topological Machine Learning} (TML) bridges this gap by integrating concepts from algebraic topology into ML workflows, enabling the discovery of patterns and features that are otherwise elusive. Despite this utility, much of the existing literature on topological methods in ML is highly technical, making it challenging for newcomers to grasp the direct connections to practical applications. 

In this tutorial, we introduce the fundamental concepts of topological methods to the machine learning (ML) community and a broader audience interested in integrating these novel approaches into their research. No prior knowledge of topology or ML is required. Our primary aim is to address this pressing need by providing a practical guide for non-experts looking to employ topological techniques in various ML contexts. To maintain accessibility, we will simplify the exposition and offer references to more detailed technical resources for those interested in further exploration.

In this paper, we teach two cornerstone techniques of TML: \textit{Persistent Homology} and \textit{Mapper algorithm}, and their effective utilization in ML. Persistent homology offers a robust, multi-scale analysis of topological features, allowing researchers to detect and quantify structures such as clusters, loops, and voids across different scales within the data. This capability is particularly beneficial for understanding the intricate relationships and hierarchies that may exist in complex datasets. From an ML perspective, PH offers \textit{a great feature extraction method} for complex datasets, which was impossible with other methods. In this part, we will focus on this aspect of PH, giving hands-on instructions with illustrations on deriving effective topological vectors from complex data. On the other hand, the Mapper algorithm complements this by providing a visual and interpretable summary of high-dimensional data. By constructing a summary graph that mirrors the underlying topology of data, the Mapper algorithm facilitates the exploration and interpretation of data in an intuitive and informative way. This technique is instrumental in uncovering data's geometric and topological essence, making it accessible for practical analysis. 


Throughout the paper, we provide comprehensive explanations and step-by-step implementations of these techniques, supported by case studies spanning diverse applications such as cancer diagnosis, shape recognition, genotyping, and drug discovery. We aim to equip researchers and practitioners with the necessary knowledge and tools to integrate TML techniques into their studies, thereby unlocking new avenues for discovery and innovation. By demonstrating the practical utility of persistent homology and the Mapper algorithm, we highlight their potential to reveal insights that traditional methods may overlook, ultimately advancing the field of Machine Learning. 

We note that this paper is not intended to be a survey of recent advances in topological data analysis but rather a tutorial aimed at introducing the fundamentals of the topic to ML practitioners. For comprehensive surveys in topological data analysis, refer to~\cite{hensel2021survey,pun2022persistent,chazal2021introduction,skaf2022topological}. For in-depth discussions of these topics, see the excellent textbooks on TDA and computational topology~\cite{dey2022computational,edelsbrunner2022computational}.

\subsection{Roadmap for the Tutorial}
  

We recommend reading the entire tutorial for a complete understanding, but readers may skip sections not pertinent to their needs. Here, we give a quick overview of the paper's structure.
%Begin with \Cref{fig:pipe} for a quick overview of the paper’s structure.

In~\Cref{sec:background}, we provide the essential topological background needed to follow the concepts discussed throughout the rest of the paper. We aim to introduce these topological concepts, help build an intuition for TDA approaches, and demonstrate how they can be adapted and applied to various needs. For those unfamiliar with topology, we strongly encourage reading our introductory crash course in~\Cref{sec:topology}. When discussing homology, we start with a non-technical, brief overview~(\Cref{sec:TLDR}), which should suffice for following the rest of the paper. For readers interested in more technical details of homology computation, we offer a more in-depth explanation in~\Cref{sec:homology2}.

 

In~\Cref{sec:PH}, we introduce \textit{Persistent Homology} (PH) in three key sections: constructing filtrations (\Cref{sec:filtrations}), deriving persistence diagrams (PD) (\Cref{sec:PD}), and applying PDs to ML tasks, including vectorization (\Cref{sec:vectorization}) and neural networks (\Cref{sec:vecNN}). The filtration process is tailored to different data formats, as methods differ significantly based on the type of data—whether it’s point clouds (\Cref{sec:point_cloud}), images (\Cref{sec:image}), or networks (\Cref{sec:graph}). If you focus on a specific data type, you can skip sections on other formats. In each subsection, we also discuss the hyperparameter selection process. Lastly, \Cref{sec:PH_software} covers available software for PH. %\Cref{fig:pipe} summarizes our PH pipeline.

In \Cref{sec:MP}, we give a brief introduction to a niche subfield, \textit{Multiparameter Persistence} (MP), an effective extension of PH. Again, we outline how to tailor the method to specific data formats for multifiltrations, i.e., point clouds (\Cref{sec:MP-point_cloud}), images~(\Cref{sec:MP-image}), and networks~(\Cref{sec:MP-graph}). Next, we outline the state-of-the-art methods on how to integrate MP information in ML pipelines~(\Cref{sec:MP-vec}).

In \Cref{sec:mapper}, we begin with a friendly introduction to the Mapper method~(\Cref{sec:mapper-PC}), an effective TML tool for unsupervised learning. We then cover hyperparameter selection for Mapper in \Cref{sec:mapperParams}, which is key for applications. Although the original Mapper algorithm is intended for point clouds, \Cref{sec:mapper-graph} explores recent advancements that extend its application to images and networks. Lastly, we provide an overview of available software libraries for Mapper in \Cref{sec:mapper-software}.

In \Cref{sec:applications}, we summarize five real-life applications of these methods from published works, namely shape recognition for point clouds (\Cref{sec:shape}), anomaly forecasting for transaction networks~(\Cref{sec:crypto}), cancer detection from histopathological images~(\Cref{sec:histo}), computer-aided drug discovery~(\Cref{sec:drug}), and cancer genotyping from RNA sequencing~(\Cref{sec:mapper-app}).

Finally, in~\Cref{sec:future}, we outline potential future directions to advance TML methods, aiming to improve their practical use in ML and discuss strategies for broadening the application of topological methods to new and emerging fields.


 